\Needspace{9\baselineskip}\begin{center}\itshape \textit{Beweis des in Abschnitt IV erwähnten Satzes.}\end{center}


\Needspace{6\baselineskip}\Article{357}
Die Gleichung vom Grad \(m\), deren Wurzeln die \(m\) Wurzeln in der Periode \((m,1)\) sind, sei \[x^m-ax^{m-1}+bx^{m-2}-\text{usw.}=0\] oder \(z=0\). Dann ist \(a=(m,1)\), und jeder der übrigen Koeffizienten \(b\) usw. hat eine Form wie \(\mathfrak A+\mathfrak B(m,1)+\mathfrak C(m,g)\), wobei \(\mathfrak A,\allowbreak \mathfrak B,\allowbreak \mathfrak C\) ganze Zahlen sind (Artikel 348). Bezeichnen wir mit \(z'\) die Funktion, in die \(z\) übergeht, wenn überall \((m,1)\) durch \((m,g)\) und \((m,g)\) durch \((m,gg)\) oder, was dasselbe ist, durch \((m,1)\) ersetzt wird, so sind die Wurzeln der Gleichung \(z'=0\) die in \((m,g)\) enthaltenen Wurzeln, und das Produkt ist \[zz'=\frac{x^n-1}{x-1}=X\]

Damit lässt sich \(z\) auf eine Form wie \(R+S(m,1)+T(m,g)\) bringen, wobei \(R,\allowbreak S,\allowbreak T\) Polynome in \(x\) sind, deren sämtliche Koeffizienten ebenfalls ganze Zahlen sind. Dann gilt \[z'=R+S(m,g)+T(m,1)\]

Schreiben wir nun der Kürze wegen \(p\) und \(q\) für \((m,1)\) bzw. \((m,g)\), so folgt \[\begin{aligned}2z&=2R+(S+T)(p+q)-(T-S)(p-q)\\&=2R-S-T-(T-S)(p-q)\end{aligned}\] und ebenso \[2z'=2R-S-T+(T-S)(p-q)\] Setzen wir daher \[2R-S-T=Y,\quad T-S=Z\]
\SourceLeaf{449}
so ergibt sich \(4X=YY-(p-q)^2ZZ\) und, da \((p-q)^2=\pm n\) gilt, \[4X=YY\mp nZZ\] Dabei gilt das obere Vorzeichen, wenn \(n\) von der Form \(4k+1\) ist, das untere, wenn \(n\) von der Form \(4k+3\) ist. Dies ist der Satz, dessen Beweis wir oben (Artikel 124) versprochen haben. Man sieht leicht, dass die beiden höchsten Glieder der Funktion \(Y\) stets \(2x^m+x^{m-1}\) sind und das höchste Glied der Funktion \(Z\) gleich \(x^{m-1}\) ist. Die übrigen Koeffizienten, die offenbar sämtlich ganze Zahlen sind, ändern sich mit der unterschiedlichen Beschaffenheit der Zahl \(n\) und lassen sich nicht durch eine allgemeine analytische Formel angeben.

\textit{Beispiel.} Für \(n=17\) ergibt sich die Gleichung, deren Wurzeln die acht in \((8,1)\) enthaltenen Wurzeln sind, nach den Regeln von Artikel 348 als \[\begin{aligned}x^8-px^7+(4+p+2q)x^6&-(4p+3q)x^5+(6+3p+5q)x^4\\&-(4p+3q)x^3+(4+p+2q)xx-px+1=0\end{aligned}\] Daraus folgt \[\begin{aligned}R&=x^8+4x^6+6x^4+4xx+1\\S&=-x^7+x^6-4x^5+3x^4-4x^3+xx-x\\T&=2x^6-3x^5+5x^4-3x^3+2xx\end{aligned}\] und hieraus \[\begin{aligned}Y&=2x^8+x^7+5x^6+7x^5+4x^4\\&\qquad+7x^3+5xx+x+2\\Z&=x^7+x^6+x^5+2x^4+x^3+xx+x\end{aligned}\] Hier sind noch einige weitere Beispiele: {\small\setlength{\arraycolsep}{4pt}\[\begin{array}{c|l|l}n&Y&Z\\\hline 3&2x+1&1\\5&2xx+x+2&x\\7&2x^3+xx-x-2&xx+x.\\11&2x^5+x^4-2x^3+2xx-x-2&x^4+x\\13&2x^6+x^5+4x^4-x^3+4xx+x+2&x^5+x^3+x\\19&2x^9+x^8-4x^7+3x^6+5x^5-5x^4&x^8-x^6+x^5+x^4-x^3+x\\{}&{}-3x^3+4xx-x-2&{}\\23&2x^{11}+x^{10}-5x^9-8x^8-7x^7-4x^6&x^{10}+x^9-x^7-2x^6-2x^5\\{}&{}+4x^5+7x^4+8x^3+5xx-x-2&{}-x^4+xx+x\end{array}\]}
\SourceLeaf{450}

\Needspace{9\baselineskip}\begin{center}\itshape \textit{Über die Gleichung für die Verteilung der Wurzeln \(\Omega\) in drei Perioden.}\end{center}


\Needspace{6\baselineskip}\Article{358}
Wir gehen zur Betrachtung der kubischen Gleichungen über, durch die für den Fall, dass \(n\) von der Form \(3k+1\) ist, die drei Summen mit \(\frac13(n-1)\) Gliedern bestimmt werden, welche die Menge \(\Omega\) bilden. Es sei \(g\) eine beliebige primitive Wurzel für den Modul \(n\) und wir setzen \(\frac13(n-1)=m\), wobei der so bestimmte Wert eine gerade ganze Zahl ist. Dann sind die drei Summen, aus denen \(\Omega\) besteht, \((m,1),\allowbreak (m,g),\allowbreak (m,gg)\), für die wir jeweils \(p,\allowbreak p',\allowbreak p''\) schreiben werden. Offenbar enthält die erste die Wurzeln \([1],\allowbreak [g^3],\allowbreak [g^6]\ldots[g^{n-4}]\), die zweite \([g],\allowbreak [g^4]\ldots[g^{n-3}]\) und die dritte \([gg],\allowbreak [g^5]\ldots[g^{n-2}]\). Setzen wir voraus, dass die gesuchte Gleichung \[x^3-Axx+Bx-C=0\] lautet, so gilt \[\begin{gathered}A=p+p'+p'',\quad B=pp'+p'p''+pp'',\\C=pp'p''\end{gathered}\]

woraus sich sogleich \(A=-1\) ergibt. Die kleinsten positiven Reste der Zahlen \(g^3\), \(g^6\ldots g^{n-4}\) nach dem Modul \(n\) seien in beliebiger Reihenfolge \(\mathfrak A,\allowbreak \mathfrak B,\allowbreak \mathfrak C\) usw., und \(\mathfrak K\) sei deren Menge nach Hinzunahme der Zahl \(1\). Ebenso seien \(\mathfrak A',\allowbreak \mathfrak B',\allowbreak \mathfrak C'\) usw. die kleinsten Reste der Zahlen \(g,\allowbreak g^4,\allowbreak g^7\ldots g^{n-3}\) und \(\mathfrak K'\) deren Menge; schließlich seien \(\mathfrak A'',\allowbreak \mathfrak B'',\allowbreak \mathfrak C''\) usw. die kleinsten Reste von \(gg,\allowbreak g^5,\allowbreak g^8\ldots g^{n-2}\) und \(\mathfrak K''\) deren Menge. Dann sind sämtliche Zahlen in \(\mathfrak K,\allowbreak \mathfrak K',\allowbreak \mathfrak K''\) verschieden und stimmen mit \(1,\allowbreak 2,\allowbreak 3\ldots n-1\) überein. Vor allem ist hier zu bemerken, dass die Zahl \(n-1\) notwendig in \(\mathfrak K\) vorkommt, da man leicht sieht, dass sie der Rest von \(g^{\frac{3m}{2}}\) ist. Daraus folgt auch leicht, dass zwei Zahlen wie \(h,\allowbreak n-h\) stets in \textit{derselben} der drei Mengen \(\mathfrak K,\allowbreak \mathfrak K',\allowbreak \mathfrak K''\) vorkommen. Ist nämlich die eine der Rest der Potenz \(g^\lambda\), so ist die andere der Rest der Potenz \(g^{\lambda+\frac{3m}{2}}\) oder der Potenz \(g^{\lambda-\frac{3m}{2}}\), wenn \(\lambda>\frac{3m}{2}\) gilt. Mit dem Zeichen \((\mathfrak K\mathfrak K)\) bezeichnen wir die Anzahl der Zahlen in der Reihe \(1,\allowbreak 2,\allowbreak 3\ldots n-1\), die sowohl selbst als auch zusammen mit den unmittelbar folgenden, um eine Einheit größeren Zahlen in \(\mathfrak K\) enthalten sind. Ebenso sei \((\mathfrak K\mathfrak K')\) die Anzahl der Zahlen in derselben Reihe, die selbst in \(\mathfrak K\), deren unmittelbare Nachfolger aber in \(\mathfrak K'\) enthalten sind. Damit ist zugleich die Bedeutung der Zeichen \((\mathfrak K\mathfrak K''),\allowbreak (\mathfrak K'\mathfrak K),\allowbreak (\mathfrak K'\mathfrak K'),\allowbreak (\mathfrak K'\mathfrak K''),\allowbreak (\mathfrak K''\mathfrak K),\allowbreak (\mathfrak K''\mathfrak K'),\allowbreak (\mathfrak K''\mathfrak K'')\) unmittelbar klar. Ich behaupte nun \textit{erstens}, dass \((\mathfrak K\mathfrak K')=(\mathfrak K'\mathfrak K)\) gilt. Sind nämlich \(h,\allowbreak h',\allowbreak h''\) usw. sämtliche Zahlen der Reihe \(1,\allowbreak 2,\allowbreak 3\ldots n-1\), die selbst in \(\mathfrak K\) enthalten sind, deren nächstgrößere Zahlen \(h+1,\allowbreak h'+1,\allowbreak h''+1\) usw. aber in \(\mathfrak K'\) enthalten sind, und deren Anzahl daher \(=(\mathfrak K\mathfrak K')\) ist, so sind offenbar sämtliche Zahlen \(n-h-1\), \(n-h'-1,\allowbreak n-h''-1\) usw. in \(\mathfrak K'\) und die nächstgrößeren \(n-h,\allowbreak n-h'\) usw.
\SourceLeaf{451}
in \(\mathfrak K\) enthalten. Da es insgesamt \((\mathfrak K'\mathfrak K)\) solche Zahlen gibt, kann gewiss nicht \((\mathfrak K'\mathfrak K)<(\mathfrak K\mathfrak K')\) gelten; ebenso beweist man, dass \((\mathfrak K\mathfrak K')<(\mathfrak K'\mathfrak K)\) nicht gelten kann. Daher müssen diese Zahlen notwendig gleich sein. Auf ganz dieselbe Weise beweist man \((\mathfrak K\mathfrak K'')=(\mathfrak K''\mathfrak K)\), \((\mathfrak K'\mathfrak K'')=(\mathfrak K''\mathfrak K')\). Zweitens: Da auf jede Zahl aus \(\mathfrak K\) mit Ausnahme der größten, \(n-1\), notwendig eine unmittelbar folgende Zahl kommt, die entweder in \(\mathfrak K\), in \(\mathfrak K'\) oder in \(\mathfrak K''\) enthalten ist, ist die Summe \((\mathfrak K\mathfrak K)+(\mathfrak K\mathfrak K')+(\mathfrak K\mathfrak K'')\) gleich der um eine Einheit verminderten Anzahl sämtlicher Zahlen in \(\mathfrak K\), also \(=m-1\). Aus einem ähnlichen Grund gilt \[\begin{aligned}(\mathfrak K'\mathfrak K)+(\mathfrak K'\mathfrak K')+(\mathfrak K'\mathfrak K'')\\{}=(\mathfrak K''\mathfrak K)+(\mathfrak K''\mathfrak K')+(\mathfrak K''\mathfrak K'')=m\end{aligned}\]

Nach diesen Vorbereitungen entwickeln wir nach den Regeln von Artikel 345 das Produkt \(pp'\) in \((m,\mathfrak A'+1)+(m,\mathfrak B'+1)+(m,\mathfrak C'+1)+\text{usw.}\). Dieser Ausdruck lässt sich, wie man leicht sieht, auf \((\mathfrak K'\mathfrak K)p+(\mathfrak K'\mathfrak K')p'+(\mathfrak K'\mathfrak K'')p''\) zurückführen. Da nach Artikel 345 I das Produkt \(p'p''\) daraus entsteht, wenn man \((m,1),\allowbreak (m,g),\allowbreak (m,gg)\) jeweils durch \((m,g),\allowbreak (m,gg)\), \((mg^3)\), das heißt \(p,\allowbreak p',\allowbreak p''\) jeweils durch \(p',\allowbreak p'',\allowbreak p\) ersetzt, gilt \(p'p''=(\mathfrak K'\mathfrak K)p'+(\mathfrak K'\mathfrak K')p''+(\mathfrak K'\mathfrak K'')p\) und auf ganz ähnliche Weise \(p''p=(\mathfrak K'\mathfrak K)p''+(\mathfrak K'\mathfrak K')p+(\mathfrak K'\mathfrak K'')p'\). Hieraus folgt unmittelbar \textit{erstens} \[B=m(p+p'+p'')=-m\] \textit{zweitens}: Da sich auf ähnliche Weise wie zuvor \(pp'\) auch \(pp''\) auf \((\mathfrak K''\mathfrak K)p+(\mathfrak K''\mathfrak K')p'+(\mathfrak K''\mathfrak K'')p''\) zurückführen lässt und dieser Ausdruck mit dem vorigen identisch sein muss, gilt notwendig \((\mathfrak K''\mathfrak K)=(\mathfrak K'\mathfrak K')\) und \((\mathfrak K''\mathfrak K'')=(\mathfrak K'\mathfrak K)\). Setzen wir nun \[\begin{aligned}(\mathfrak K'\mathfrak K'')&=(\mathfrak K''\mathfrak K')=a,\\(\mathfrak K''\mathfrak K'')&=(\mathfrak K'\mathfrak K)=(\mathfrak K\mathfrak K')=b,\\(\mathfrak K'\mathfrak K')&=(\mathfrak K''\mathfrak K)=(\mathfrak K\mathfrak K'')=c\end{aligned}\] so ergibt sich \(m-1=(\mathfrak K\mathfrak K)+(\mathfrak K\mathfrak K')+(\mathfrak K\mathfrak K'')=(\mathfrak K\mathfrak K)+b+c\) und \(a+b+c=m\), also \((\mathfrak K\mathfrak K)=a-1\). Damit sind jene neun unbekannten Größen auf drei, \(a,\allowbreak b,\allowbreak c\), oder vielmehr wegen der Gleichung \(a+b+c=m\) auf zwei zurückgeführt. Schließlich ist klar, dass sich das Quadrat \(pp\) in \((m,1+1)+(m,\mathfrak A+1)+(m,\mathfrak B+1)+(m,\mathfrak C+1)+\text{usw.}\) entwickeln lässt. Unter den Teilen dieses Ausdrucks findet sich \((m,n)\), das auf \((m,0)\) oder auf \(m\) zurückgeführt wird; die übrigen Teile lassen sich offenbar auf \((\mathfrak K\mathfrak K)p+(\mathfrak K\mathfrak K')p'+(\mathfrak K\mathfrak K'')p''\) zurückführen, sodass \(pp=m+(a-1)p+bp'+cp''\) gilt.

So haben wir durch die vorstehenden Untersuchungen diese vier Umformungen erhalten.
\SourceLeaf{452}\[\begin{aligned}pp&=m+(a-1)p+bp'+cp''\\pp'&=bp+cp'+ap''\\pp''&=cp+ap'+bp''\\p'p''&=ap+bp'+cp''\end{aligned}\] Dabei besteht zwischen den drei Unbekannten \(a,\allowbreak b,\allowbreak c\) die Bedingungsgleichung \[a+b+c=m\quad \ldotp\,\ldotp\,\ldotp\,\ldotp\,\ldotp\,\ldotp\,\ldotp\,\ldotp\,\ldotp\,\ldotp\,\ldotp\ (\mathrm{I})\] und außerdem ist sicher, dass sie ganze Zahlen sind. Daraus folgt \[\begin{aligned}C=p\times p'p''&=app+bpp'+cpp''\\&=am+(aa+bb+cc-a)p\\&\quad+(ab+bc+ac)p'+(ab+bc+ac)p''\end{aligned}\]

Da aber \(pp'p''\) eine symmetrische Funktion der Summen \(p,\allowbreak p',\allowbreak p''\) ist, müssen die Koeffizienten, mit denen diese im vorstehenden Ausdruck multipliziert werden, notwendig gleich sein (Artikel 350). Daher erhalten wir die neue Gleichung \[aa+bb+cc-a=ab+bc+ac\quad \ldotp\,\ldotp\,\ldotp\ (\mathrm{II})\] und daraus \(C=am+(ab+bc+ac)(p+p'+p'')\) oder, wegen I und \(p+p'+p''=-1\), \[C=aa-bc\quad \ldotp\,\ldotp\,\ldotp\,\ldotp\,\ldotp\,\ldotp\ (\mathrm{III})\]

Obwohl nun \(C\) hier von drei Unbekannten abhängt, zwischen denen nur zwei Gleichungen bestehen, genügen diese dennoch zusammen mit der Bedingung, dass \(a,\allowbreak b,\allowbreak c\) ganze Zahlen sind, zur vollständigen Bestimmung von \(C\). Um dies zu zeigen, schreiben wir die Gleichung II als \[\begin{aligned}12a+12b+12c+4&=36aa+36bb+36cc\\&\quad-36ab-36ac-36bc\\&\quad-24a+12b+12c+4\end{aligned}\] Die linke Seite wird nach I zu \(=12m+4=4n\); die rechte lässt sich auf \[(6a-3b-3c-2)^2+27(b-c)^2\] oder, wenn wir \(k\) für \(2a-b-c\) schreiben, auf \((3k-2)^2+27(b-c)^2\) zurückführen. Daraus ist klar, dass sich die Zahl \(4n\) (das heißt allgemein das Vierfache einer beliebigen Primzahl von der Form \(3m+1\)) durch die Form \(xx+27yy\) darstellen lässt. Dies lässt sich zwar ohne Schwierigkeit aus der allgemeinen Theorie
\SourceLeaf{453}
der binären Formen ableiten, doch ist es recht bemerkenswert, dass eine solche Zerlegung mit den Werten von \(a,\allowbreak b,\allowbreak c\) zusammenhängt. Die Zahl \(4n\) lässt sich aber stets nur auf eine einzige Weise in ein Quadrat und das \(27\)-Fache eines Quadrats zerlegen, was wir folgendermaßen beweisen\HistoricalNote{*)}{Dieser Satz ließe sich unmittelbarer aus den Grundsätzen von Abschnitt V beweisen.}. Setzen wir \[4n=tt+27uu=t't'+27u'u'\] voraus, so gilt \textit{erstens} \[(tt'-27uu')^2+27(tu'+t'u)^2=16nn\] \textit{zweitens} \[(tt'+27uu')^2+27(tu'-t'u)^2=16nn\] \textit{drittens} \[(tu'+t'u)(tu'-t'u)=4n(u'u'-uu)\] Aus der dritten Gleichung folgt, dass \(n\), da es eine Primzahl ist, eine der Zahlen \(tu'+t'u,\allowbreak tu'-t'u\) teilt; aus der ersten und zweiten ist aber klar, dass jede dieser Zahlen kleiner als \(n\) ist. Daher muss diejenige, die durch \(n\) teilbar ist, notwendig \(=0\) sein, und somit auch \(u'u'-uu=0\), woraus \(u'u'=uu\) und \(t't'=tt\) folgt; jene beiden Zerlegungen unterscheiden sich also nicht. Setzen wir daher die Zerlegung von \(4n\) in ein Quadrat und das \(27\)-Fache eines Quadrats als bekannt voraus (man kann sie entweder nach dem direkten Verfahren von Abschnitt V oder nach dem in den Artikeln 323 und 324 angegebenen indirekten Verfahren ermitteln), gilt also \(4n=MM+27NN\), so sind die Quadrate \((3k-2)^2,\allowbreak (b-c)^2\) bestimmt, und wir haben nun anstelle der Gleichung II zwei Gleichungen erhalten. Man sieht aber leicht, dass nicht nur das Quadrat \((3k-2)^2\), sondern auch die Wurzel \(3k-2\) selbst vollständig bestimmt ist. Da diese notwendig entweder \(=+M\) oder \(=-M\) sein muss, wird die Zweideutigkeit dadurch beseitigt, dass \(k\) eine ganze Zahl sein muss. Daher setzt man \(3k-2=+M\) oder \(=-M\), je nachdem, ob \(M\) von der Form \(3z+1\) oder \(3z+2\) ist\HistoricalNote{\(\dagger\))}{Offenbar kann \(M\) nicht von der Form \(3z\) sein, da sonst \(4n\) durch \(3\) teilbar wäre. — Die Zweideutigkeit, ob \(b-c\) \(=N\) oder \(=-N\) zu setzen ist, braucht hier nicht berücksichtigt zu werden und lässt sich auch der Natur der Sache nach auf keine Weise beseitigen, da sie von der Wahl der primitiven Wurzel \(g\) abhängt: Für einige primitive Wurzeln ist die Differenz \(b-c\) positiv, für andere negativ.}. Da nun \(k=2a-b-c=3a-m\) gilt, ist \(a=\frac13(m+k),\allowbreak \ b+c=m-a=\frac13(2m-k)\), sodass \[\begin{aligned}C=aa-bc&=aa-\tfrac14(b+c)^2+\tfrac14(b-c)^2\\&=\tfrac19(m+k)^2-\tfrac1{36}(2m-k)^2+\tfrac14NN\\&=\tfrac1{12}kk+\tfrac13km+\tfrac14NN\end{aligned}\] gilt; damit sind sämtliche Koeffizienten der gesuchten Gleichung gefunden. \(Q.E.F.\) — Diese Formel
\SourceLeaf{454}
wird noch einfacher, wenn für \(NN\) sein Wert aus der Gleichung \((3k-2)^2+27NN=4n=12m+4\) eingesetzt wird; nach ausgeführter Rechnung ergibt sich \[C=\tfrac19(m+k+3km)=\tfrac19(m+kn)\] Derselbe Wert lässt sich auch auf \((3k-2)NN+k^3-2kk+k-km+m\) zurückführen. Dieser Ausdruck ist zwar zur Anwendung weniger geeignet, zeigt aber unmittelbar, dass \(C\), wie es sein soll, sicher eine ganze Zahl ist.

\textit{Beispiel.} Für \(n=19\) gilt \(4n=49+27\), also \(3k-2=+7,\allowbreak k=3\), \(C=\frac19(6+57)=7\), und die gesuchte Gleichung ist \(x^3+xx-6x-7=0\), wie oben (Artikel 351). — Ebenso ergibt sich für \(n=7,\allowbreak 13,\allowbreak 31,\allowbreak 37,\allowbreak 43,\allowbreak 61,\allowbreak 67\) der Wert von \(k\) jeweils als \(1,\allowbreak -1,\allowbreak 2,\allowbreak -3,\allowbreak -2,\allowbreak 1,\allowbreak -1\), woraus \(C=1,\allowbreak -1,\allowbreak 8,\allowbreak -11,\allowbreak -8,\allowbreak 9,\allowbreak -5\) folgt.

Obwohl das in diesem Artikel gelöste Problem recht verwickelt ist, wollten wir es dennoch nicht unterdrücken, sowohl wegen der Eleganz seiner Lösung als auch, weil es Gelegenheit gab, verschiedene Kunstgriffe anzuwenden, die auch bei anderen Fragen mit großem Nutzen verwendet werden können\HistoricalNote{*)}{Folgerung. Es sei \(\varepsilon\) eine Wurzel der Gleichung \(x^3-1=0\); dann gilt \((p+\varepsilon p'+\varepsilon\varepsilon p'')^3=\frac n2(M+N\surd -27)\). Man setze \(\dfrac M{\surd 4n}=\cos\varphi,\allowbreak \quad\dfrac{N\surd 27}{\surd 4n}=\sin\varphi\); dann gilt \[\begin{gathered}p=-\tfrac13+\tfrac23\cos\tfrac13\varphi\surd n;\\M\equiv+1\ (\operatorname{mod.}3);\quad1\equiv M(1.2.3\ldots m)^3\ (\operatorname{mod.}n)\end{gathered}\] Setzt man \(3x+1=y\) so wird die Gleichung \(y^3-3ny-Mn=0\).}.

